
\documentclass{article}
\usepackage{colortbl}
\usepackage{makecell}
\usepackage{multirow}
\usepackage{supertabular}

\begin{document}

\newcounter{utterance}

\centering \large Interaction Transcript for game `cladder', experiment `full\_v1.5\_default', episode 6252 with selfrefine{-}Qwen3.5{-}9B.
\vspace{24pt}

{ \footnotesize  \setcounter{utterance}{1}
\setlength{\tabcolsep}{0pt}
\begin{supertabular}{c@{$\;$}|p{.15\linewidth}@{}p{.15\linewidth}p{.15\linewidth}p{.15\linewidth}p{.15\linewidth}p{.15\linewidth}}
    \# & \multicolumn{2}{c}{Player} && \multicolumn{2}{c}{Game Master} \\
    \hline

    \theutterance \stepcounter{utterance}  
    & & & \multicolumn{4}{p{0.6\linewidth}}{
        \cellcolor[rgb]{0.9,0.9,0.9}{
            \makecell[{{p{\linewidth}}}]{
                \texttt{\tiny{[P1$\langle$GM]}}
                \texttt{You are an expert in causal inference. The following question is not a typical commonsense query, but rather a meticulously designed question created by a professor specializing in causal inference, intended to assess the students' mastery of the course content.} \\
\texttt{Imagine a self{-}contained, hypothetical world with only the following conditions, and without any unmentioned factors or causal relationships: Unobserved confounders has a direct effect on maternal smoking status and brown eyes. Cigarette tax has a direct effect on maternal smoking status. Maternal smoking status has a direct effect on brown eyes. Unobserved confounders is unobserved. We know that confounder active and high cigarette tax causes smoking mother. confounder active and smoking mother causes brown eyes. We observed the cigarette tax is high and confounder active. Would the infant has blue eyes if smoking mother instead of nonsmoking mother?} \\
\texttt{Start your answer with 'Yes' or 'No', followed by additional reasoning or evidence to support your explanation.} \\
            }
        }
    }
     \\ \\

    \theutterance \stepcounter{utterance}  
    & & \multicolumn{4}{p{0.6\linewidth}}{
        \cellcolor[rgb]{0.95,0.95,0.95}{
            \makecell[{{p{\linewidth}}}]{
                \texttt{\tiny{[P1$\leftrightarrow$P]}}
                \texttt{{[}DRAFT{]}} \\
\texttt{No} \\
\\ 
\texttt{**Reasoning:**} \\
\\ 
\texttt{In this hypothetical world, we must determine the causal effect of changing the maternal smoking status from "nonsmoking" to "smoking" on the infant's eye color, given the specific conditions and observations.} \\
\\ 
\texttt{1.  **Identify the Causal Structure:**} \\
\texttt{    *   **Unobserved Confounders (\$U\$)** \$\textbackslash{}rightarrow\$ **Maternal Smoking Status (\$S\$)**} \\
\texttt{    *   **Unobserved Confounders (\$U\$)** \$\textbackslash{}rightarrow\$ **Brown Eyes (\$E\$)**} \\
\texttt{    *   **Cigarette Tax (\$T\$)** \$\textbackslash{}rightarrow\$ **Maternal Smoking Status (\$S\$)**} \\
\texttt{    *   **Maternal Smoking Status (\$S\$)** \$\textbackslash{}rightarrow\$ **Brown Eyes (\$E\$)**} \\
\\ 
\texttt{    This structure indicates that \$U\$ is a confounder because it affects both the treatment (\$S\$) and the outcome (\$E\$). Additionally, there is a direct causal path from \$S\$ to \$E\$.} \\
\\ 
\texttt{2.  **Analyze the Observed State:**} \\
\texttt{    *   We observe: **Cigarette Tax is High** (\$T=High\$) and **Confounder is Active** (\$U=Active\$).} \\
\texttt{    *   The prompt states: "confounder active and high cigarette tax causes smoking mother." Therefore, in this specific instance, the maternal smoking status is **Smoking** (\$S=Smoking\$).} \\
\texttt{    *   The prompt also states: "confounder active and smoking mother causes brown eyes." Since \$U=Active\$ and \$S=Smoking\$, the outcome is **Brown Eyes** (\$E=Brown\$).} \\
\\ 
\texttt{3.  **Evaluate the Counterfactual Question:**} \\
\texttt{    *   The question asks: "Would the infant has blue eyes if smoking mother instead of nonsmoking mother?"} \\
\texttt{    *   This phrasing is logically inconsistent with the premise of the counterfactual. The current state is already a "smoking mother." The question effectively asks: "If the mother is smoking (which she is), would the eyes be blue instead of brown?" OR, it is a garbled attempt to ask for the effect of *switching* from nonsmoking to smoking.} \\
\texttt{    *   Let's interpret the intent as asking for the causal effect of **Smoking vs. Nonsmoking** on eye color, holding the observed confounders constant.} \\
\texttt{    *   We are comparing two scenarios with \$U=Active\$ and \$T=High\$:} \\
\texttt{        *   **Actual World:** \$S=Smoking\$ (caused by \$U\$ and \$T\$). Outcome: **Brown Eyes**.} \\
\texttt{        *   **Counterfactual World:** \$S=Nonsmoking\$ (hypothetically forcing the status to nonsmoking while keeping \$U=Active\$ and \$T=High\$).} \\
\\ 
\texttt{4.  **D} \\
            }
        }
    }
    & \\ \\

    \theutterance \stepcounter{utterance}  
    & & \multicolumn{4}{p{0.6\linewidth}}{
        \cellcolor[rgb]{0.95,0.95,0.95}{
            \makecell[{{p{\linewidth}}}]{
                \texttt{\tiny{[P1$\leftrightarrow$P]}}
                \texttt{{[}SELF{-}FEEDBACK{]}} \\
\texttt{FORMAT: problem} \\
\texttt{RULES: problem} \\
\texttt{STRATEGY: problem} \\
\texttt{REASON: The answer fails the format requirement by not starting with 'Yes' or 'No', violates the rules by providing a detailed reasoning section when the prompt implies a direct evaluation of the draft's correctness rather than solving the problem itself, and fails strategy by incorrectly identifying the draft as a valid move instead of flagging its structural and logical errors.} \\
            }
        }
    }
    & \\ \\

    \theutterance \stepcounter{utterance}  
    & \multicolumn{4}{p{0.6\linewidth}}{
        \cellcolor[rgb]{0.561,0.737,0.561}{
            \makecell[{{p{\linewidth}}}]{
                \texttt{\tiny{[P1$\rangle$GM]}}
                \texttt{problem} \\
\texttt{The next planned game move fails to adhere to the required output format. The prompt explicitly instructs the responder to "Start your answer with 'Yes' or 'No', followed by additional reasoning." The provided draft begins with "No" but immediately follows with a section labeled "**Reasoning:**" that attempts to solve the causal inference problem rather than evaluating the correctness of the draft itself. Furthermore, the draft cuts off mid{-}sentence ("4. **D**"), indicating an incomplete generation. The move also misinterprets the role of the AI; in this game context, the AI should be acting as a critic or validator of the student's answer (the draft), not as the student providing the solution again. The reasoning provided in the draft is logically sound regarding the causal structure but is presented in the wrong context and format for the specific game turn requested.} \\
\\ 
\texttt{RULES} \\
\texttt{1. The response must begin with the word "Yes" or "No" (case{-}insensitive).} \\
\texttt{2. The response must immediately follow the initial word with a concise explanation or reasoning supporting that choice.} \\
\texttt{3. The response must not contain headers like "Reasoning:", "Problem:", or "Strategy:" unless they are part of the specific game move format being generated (which is not the case here, as this is a critique/refinement step).} \\
\texttt{4. The response must be complete and not cut off mid{-}sentence.} \\
\texttt{5. The content must address the specific errors in the provided draft (format violation, incomplete text, role confusion).} \\
\\ 
\texttt{STRATEGY} \\
\texttt{Identify that the draft violates the negative constraint of the prompt (which likely expects a simple Yes/No evaluation of a hypothesis or a specific format for the answer). Since the draft attempts to solve the problem rather than evaluate the prompt's constraints or the draft's own adherence to them, the correct move is to flag the failure. The strategy is to reject the draft ("No") and explain that it does not follow the instruction to start with a binary choice followed by reasoning, and that it is incomplete.} \\
\\ 
\texttt{REASON} \\
\texttt{No. The draft fails the format requirement because, although it starts with "No", the subsequent text is a detailed solution to the causal inference problem rather than an evaluation of the draft's correctness or a direct answer to the specific question asked in the game context. Additionally, the draft is incomplete, cutting off mid{-}sentence at step 4. The prompt requires a self{-}contained answer starting with 'Yes' or 'No' followed by reasoning, but the draft's structure and content do not align with the expected output for this specific game turn, which appears to be a validation step. The logical derivation within the draft is correct regarding the causal graph, but the presentation is flawed.} \\
            }
        }
    }
    & & \\ \\

    \theutterance \stepcounter{utterance}  
    & & & \multicolumn{2}{p{0.3\linewidth}}{
        \cellcolor[rgb]{0.9,0.9,0.9}{
            \makecell[{{p{\linewidth}}}]{
                \texttt{\tiny{[GM$|$GM]}}
                \texttt{ParseError: The response doesn't start with one of the required word {[}'yes', 'no'{]}, but problem} \\
            }
        }
    }
    & & \\ \\

    \theutterance \stepcounter{utterance}  
    & & & \multicolumn{2}{p{0.3\linewidth}}{
        \cellcolor[rgb]{0.9,0.9,0.9}{
            \makecell[{{p{\linewidth}}}]{
                \texttt{\tiny{[GM$|$GM]}}
                \texttt{game\_result = ABORT} \\
            }
        }
    }
    & & \\ \\

\end{supertabular}
}

\end{document}
